% =============================================================
%  Fundamentação Teórica (2,50 pontos)
%  As citações \cite{...} referenciam chaves definidas em
%  referencias.bib. Adicione novas fontes lá.
% =============================================================
\section*{Fundamentação Teórica}

A literatura relevante a este projeto se organiza em quatro frentes:

\textbf{O trabalho do perito de vídeo.} Laboratórios de perícia audiovisual descrevem o exame de vídeo como um processo que envolve autenticação do material, tratamento de imagem, identificação de pessoas, objetos e placas de veículos, e produção de um laudo tecnicamente fundamentado - trabalho hoje conduzido de forma majoritariamente manual, sem padronização nacional de ferramentas entre os institutos de criminalística estaduais \cite{ipog_pericia_audiovisual,mppr_pericias_audiovisuais,oliveira2022comparativo}. Este projeto pretende aprofundar esse mapeamento na revisão de literatura, detalhando as etapas, os critérios técnicos e os desafios operacionais relatados por peritos e pesquisadores da área (autenticidade, cadeia de custódia, qualidade do material, ausência de softwares padronizados).

\textbf{Algoritmos clássicos de visão computacional para vídeo de segurança.} Modelos de detecção de objetos em tempo real (família YOLO, SSD) e algoritmos de rastreamento permitem identificar e seguir pessoas, veículos e objetos ao longo de um vídeo, servindo de base para tarefas como contagem, reconhecimento de eventos e triagem de trechos relevantes em gravações longas \cite{sciencedirect_anomaly_detection,arxiv_video_analytics_deployment}. Estudos recentes também aplicam redes convolucionais (CNN 3D, arquiteturas Inception) à detecção de anomalias e comportamentos suspeitos em circuitos fechados de televisão \cite{researchgate_suspicious_activity}, e modelos baseados em ConvLSTM à detecção de manipulações (deepfakes) em vídeos de vigilância \cite{springer_deepfake_cctv} - relevante para a verificação de autenticidade do material antes de sua análise.

\textbf{Modelos multimodais de linguagem para vídeo.} Modelos como LLaVA, MiniGPT4-Video, VideoChat e arquiteturas de memória de longo prazo (MA-LMM) combinam codificadores visuais com modelos de linguagem para gerar descrições, responder perguntas e raciocinar sobre o conteúdo de vídeos em linguagem natural \cite{arxiv_minigpt4video,arxiv_malmm,arxiv_video_llm_survey}. Pesquisas recentes também avaliam esses modelos diretamente em tarefas forenses, como identificação e descrição de manipulações em imagens \cite{arxiv_gpt_forensics}, indicando potencial de uso como camada de síntese e comunicação sobre saídas de algoritmos especializados.

\textbf{Lacunas identificadas.} A literatura sobre processamento de vídeo de vigilância é extensa, mas majoritariamente voltada a segurança em tempo real (prevenção, monitoramento ativo), não ao contexto pericial pós-fato, com suas exigências específicas de rastreabilidade, precisão documental e validade probatória. Da mesma forma, a literatura sobre modelos multimodais de linguagem para vídeo raramente aborda sua combinação com algoritmos especializados de visão computacional em um fluxo orientado à produção de laudo pericial. É nessa lacuna - entre capacidade técnica dos algoritmos disponíveis e as exigências reais do trabalho pericial - que este projeto se insere.
